% ===== 摘要 =====
\begin{abstract}
混凝土骨料粒径级配直接影响混凝土工作性能与力学性能，但传统机械筛分依赖离线取样、振筛与逐级称量，难以满足生产环节的快速反馈需求。针对自然堆积场景下骨料颗粒接触粘连、二维投影尺寸与机械筛孔语义不一致以及视觉颗粒数量分布与标准筛分质量分布口径不同等问题，本文构建了一套“实例感知--物理测量--筛分语义映射--形貌/异常--自动报告”的视觉筛分方法。系统复用 ImageGrains 2.0 / Cellpose-SAM 预训练模型完成颗粒实例分割，通过现场尺度标定将像素几何转换为毫米量纲，并以筛分等效粒径 $d=\theta_1 b+\theta_2d_{\mathrm{eq}}+\theta_3$ 和质量代理权重 $w=d^\gamma$ 将逐颗粒二维测量聚合为累计质量级配、$D_{10}/D_{50}/D_{90}$ 及标准筛级质量占比；同一实例几何表征进一步用于二维投影形貌分类与尺寸/疑似泥团异常筛查。为避免把经验基线误解释为标准筛分精度，本文将尺度标定与筛分语义参数校准分离，并设计同批图像--机械筛分配对协议，以校准批次拟合 $\boldsymbol{\theta},\gamma$、以独立批次报告最终误差。当前三张真实自然堆积演示图分别得到 1924、968 和 966 个实例，数量型与 $d^3$ 质量代理型 $D_{50}$ 呈现显著差异，说明从颗粒计数到质量语义的转换不可省略；仓库竞赛层 27 项合成测试验证了等效粒径、加权统计、校准、形貌、异常与报告模块的软件一致性。当前版本尚缺与演示样本配对的逐筛机械称量，因此上述结果属于功能与模型行为验证，而非最终筛分准确率结论。本文进一步给出完整的 30~min 现场部署流程、证据边界和物理验证路线，为后续以标准筛分真值完成参数校准和独立精度评价提供可复现框架。

\noindent\textbf{关键词：} 混凝土骨料；智能筛分；实例分割；Cellpose-SAM；粒径级配；质量加权；物理校准
\end{abstract}
